\chapter{Storage}

\section{What the World Database Actually Has to Hold}

Chapter 25 left the world database as a black box: something that holds \(M_t\), accepts reads, and accepts writes only from the commit stage. It is worth being precise about what that actually means before designing anything, because \(M_t\) is not a single object of fixed size. It is five components, Φ, v, S, R, and z, each potentially defined over an unbounded domain, and Chapter 8 already established that most of that domain, at any given moment, is not densely resolved at all. A naive implementation, allocating storage for a dense array of all five components across the entire domain whether or not anything has actually been observed there, would waste an amount of space bounded only by how large the domain is permitted to grow, most of it spent representing regions that hold nothing more specific than an unresolved live possibility set. This chapter's task is to design storage that reflects the field's actual structure rather than a worst-case assumption about it.

\section{Sparse by Construction}

The entropy constraint from Chapter 8, together with the invention-to-memory discipline from Chapter 9, together imply something useful for storage design specifically: a region only acquires dense, specific, low-entropy content once it has actually been observed. Everything else remains, honestly, at high entropy, represented by \(L_t(x)\) rather than by a committed z. This is exactly the structure a sparse storage scheme is built to exploit. Rather than a dense array, the world database should use a hierarchical, sparse structure, an octree or quadtree depending on the domain's dimensionality, that allocates storage only for regions actually holding resolved content, with unresolved regions represented implicitly, by their absence from the structure, rather than explicitly, by a stored placeholder. A courtyard's unobserved far corner costs the storage system nothing beyond whatever coarse, cheap representation of \(L_t(x)\) is useful for the entropy constraint's own bookkeeping, until the moment it is actually observed and Chapter 9's invention commits real content there.

\section{Locality and Space-Filling Curves}

Sparse allocation solves the size problem. It does not, by itself, solve a related problem: a scheme that allocates storage sparsely can still scatter logically nearby points across physically distant storage locations, which is costly for exactly the kind of patch-based access Chapter 21's cover-based construction requires, reading a whole neighborhood's worth of points together to check a gluing condition or perform a projection. The standard technique for keeping spatial neighbors close in storage as well as in the field's own domain is a space-filling curve, most commonly a Hilbert curve, which linearizes a multi-dimensional domain into a one-dimensional storage order while preserving locality far better than a naive row-major traversal: two points close together in the field's domain are, with high probability, also close together in the resulting storage order. This matters concretely for this book's own machinery. Chapter 21's patches, the units over which local admissibility is checked and gluing conditions evaluated, correspond naturally to contiguous ranges along a Hilbert-ordered storage layout, so that loading a patch means reading one contiguous span of storage rather than gathering scattered fragments from across the whole structure.

Purpose-driven reordering of this kind is not a novelty invented for computing. The Arabic script offers a centuries-old precedent: alongside the ancient Abjad order, which preserves inherited numerical and mnemonic structure and remains in use for numerals and chronology, a separate Hijā'ī order was later adopted, regrouping the same twenty-eight letters by shared shape purely to ease handwriting instruction and lookup, without displacing Abjad's continued use where its own structure is still what is needed. A Hilbert-ordered storage layout is doing something with the same shape: keeping every point the field's domain contains while imposing a second, purpose-built ordering optimized for an access pattern the domain's own coordinates were never designed to serve. Chapter 28 develops the comparison in full.

\section{Entropy-Dependent Compression}

This gives entropy, Chapter 8's S, a genuine downstream engineering role beyond the constraint it was originally introduced to enforce. A region's entropy is a direct, principled signal for how much storage that region deserves: high-S regions need only enough representation to describe \(L_t(x)\) at whatever coarse resolution the entropy constraint's own bookkeeping requires, while low-S regions, fully resolved, earn full-resolution storage for their committed z and whatever residue R has accumulated there. Compression ratio, in this scheme, is not a separate parameter tuned independently of the field's own mathematics; it is a direct function of entropy, high where entropy is high and low where entropy is low, which means the storage system's compression behavior is automatically doing the right thing wherever Chapter 8's constraint is already being honestly enforced elsewhere in the system.

\section{Different Strategies for Different Components}

It is worth being explicit that the five components do not warrant identical storage treatment, because they do not share the same character. Appearance, z, is dense, image-like data wherever it is resolved, and benefits from whatever general-purpose compression suits that kind of content. Residue, R, is not a full event log, and Chapter 11 was explicit about why: it is a compressed, structural trace, and its storage representation should reflect that directly, closer to a compact causal-graph fragment than to a growing append-only journal, exactly matching the ⊕-accumulated object Chapter 11 described rather than anything larger. Φ, v, and S, given Chapter 17's fast-slow distinction, can generally tolerate coarser storage resolution and less frequent versioning than z: a coherence field that changes on a slow clock does not need the same fine-grained temporal snapshotting a fast-changing appearance field does, and storing it as though it did would waste space tracking versions of a quantity that, across most of those versions, has not actually changed.

\section{Streaming and Progressive Loading}

A world database of any real size will not fit entirely in memory at once, and the sparse, patch-organized structure this chapter has built turns out to make this manageable rather than merely tolerable. Because Chapter 21's patches are already the natural unit of both mathematical admissibility-checking and Hilbert-ordered storage locality, they are also the natural unit of loading and unloading: as an observer moves through the domain, patches near the observer's current position are streamed into memory, and patches far from any current observer are evicted, without this eviction meaning anything has been lost, since the world database, not memory, remains the authoritative store. This chapter will not develop the full distributed-systems treatment of a world database too large for any single machine; that belongs to Chapter 30. What matters here is that streaming is not a bolt-on feature requiring a separate data structure. It falls directly out of organizing storage around the same patches Chapter 21 already needed for entirely different, purely mathematical reasons.

\section{A Storage Interface}

It is worth gesturing, briefly, at what a concrete implementation of this chapter's design would actually expose to the rest of the system, without providing a full specification. A storage backend, in this architecture, needs to implement a narrow interface: given a patch identifier, retrieve its current committed content across whatever subset of Φ, v, S, R, and z the caller needs; given a patch identifier and a committed update, write it, respecting whatever compression and versioning policy section 26.4 and 26.5 established for that component; and given a region of the domain, enumerate which patches currently exist within it, distinguishing resolved patches from the implicitly-represented unresolved space around them. Everything this chapter has described, sparsity, Hilbert ordering, entropy-dependent compression, component-specific strategies, streaming, lives behind this interface rather than being exposed to the proposal, projection, and commit services from Chapter 25, each of which only ever needs to know that the world database can answer these three kinds of request, not how it answers them.

\section{Looking Ahead}

This chapter has given the world database a real storage design: sparse, entropy-aware, organized for locality, and streamable. What it has not addressed is what happens on the other side of a read, once a patch's content has actually been retrieved by a client. Chapter 27 takes up rendering directly, asking how the appearance a storage system hands back actually becomes an observation an agent can use.
